\title{QLoRA: Efficient Finetuning of Quantized LLMs}

\begin{abstract}
We introduce \textsc{QLoRA}, a memory-efficient finetuning strategy that enables the adaptation of a 65B parameter model on a single 48GB GPU, without compromising the task performance typical of full 16-bit finetuning. With \textsc{QLoRA}, gradients are propagated through a static, 4-bit quantized pretrained language model into Low Rank Adapters~(LoRA). Our top-performing model series, \textbf{Guanaco}, sets a new benchmark by surpassing all prior publicly available models on the Vicuna evaluation, attaining 99.3\% of ChatGPT's performance after only 24 hours of single-GPU finetuning. Several innovations underlie \textsc{QLoRA}'s memory efficiency and strong performance: (a) 4-bit NormalFloat (NF4), a novel datatype theoretically optimal for weights following a normal distribution, (b) Double Quantization, which minimizes memory usage further by quantizing the quantization parameters themselves, and (c) Paged Optimizers that address memory peak management. Employing \textsc{QLoRA}, we finetune over 1,000 models and systematically investigate instruction following and chatbot abilities across 8 instruction datasets, diverse model architectures (LLaMA, T5), and model sizes otherwise prohibitively large for standard finetuning practices (such as 33B and 65B parameter models). The results indicate that QLoRA finetuning on concise, high-quality datasets achieves cutting-edge outcomes, even with smaller models than those used in previous state-of-the-art systems. We present an in-depth assessment of chatbot capabilities via both human evaluators and GPT-4, illustrating that GPT-4 can serve as an inexpensive yet effective stand-in for human judgment. Our findings also reveal that current chatbot benchmarks have significant limitations for accurately measuring chatbot proficiency. Through targeted, failure-focused analysis, we show specific scenarios where \textbf{Guanaco} does not match ChatGPT's performance. All models, source code, and 4-bit CUDA training kernels are publicly released.\footnote{\url{https://github.com/artidoro/qlora} and \url{https://github.com/TimDettmers/bitsandbytes}}
\end{abstract}

\section{Introduction}

Finetuning large language models (LLMs) has proven to be an efficient strategy for boosting their capabilities \citep{min2021metaicl, wei2021finetuned, ouyang2022training, wang2022super, wang2022self, liu2022few}, as well as for integrating preferred behaviors or mitigating unwanted ones \citep{ouyang2022training,askell2021general,bai2022training}. Nevertheless, adapting very large models comes with enormous computational demands; for instance, standard 16-bit finetuning of the LLaMA model with 65B parameters~\citep{touvron2023llama} consumes upwards of 780 GB of GPU memory. Although advancements in quantization have made it feasible to lower the memory consumption of LLMs \citep{dettmers2022llm, dettmers2022case, frantar2022gptq, xiao2022smoothquant}, such methods remain restricted to inference tasks and are not applicable during training \citep{wortsman2023stable}. 

In this work, we present the first approach that enables the successful finetuning of a quantized 4-bit model while fully preserving its performance. Our proposed technique, \textsc{QLoRA}, introduces an innovative high-precision quantization method that compresses a pretrained model to 4 bits. Subsequently, it incorporates a lightweight set of trainable Low-rank Adapter weights \citep{hu2021lora}, 

which are optimized by propagating gradients through the quantized model parameters. 

\textsc{QLoRA} decreases the average memory usage for finetuning a 65B parameter model from more than 780GB to under 48GB of GPU memory, without sacrificing either training speed or predictive accuracy when compared to full 16-bit finetuning. This development marks a considerable enhancement in the accessibility of LLM finetuning, making it feasible to finetune the largest publicly available models using a single GPU. Utilizing \textsc{QLoRA}, we introduce the \textbf{Guanaco} model suite. The second-best \textbf{Guanaco} model achieves 97.8\% of ChatGPT's performance on the Vicuna~\citep{vicuna2023} benchmark, while the model can be trained in under 12 hours on a single consumer GPU; the largest model reaches 99.3\% in 24 hours with one professional GPU, effectively matching ChatGPT's performance on Vicuna. At deployment, our most compact \textbf{Guanaco} model, with 7B parameters, operates using only 5 GB of memory, surpassing the performance of the 26 GB Alpaca model by over 20 percentage points on the Vicuna benchmark (see Table~\ref{tbl:vicuna_eval}).

\textsc{QLoRA} presents several key innovations aimed at minimizing memory consumption while maintaining high performance levels: (1) \textbf{4-bit NormalFloat}, a quantization format optimized for data following a normal distribution, which achieves superior empirical outcomes compared to both 4-bit Integer and 4-bit Float representations. (2) \textbf{Double Quantization}, a technique that applies quantization to the quantization coefficients themselves, resulting in average savings of roughly 0.37 bits per parameter (translating to around 3 GB for a 65B parameter model). (3) \textbf{Paged Optimizers}, which leverage NVIDIA's unified memory framework to eliminate the memory surges associated with gradient checkpointing during the processing of mini-batches with extended sequence lengths. These advances are integrated into a refined LoRA methodology that incorporates adapters throughout all network layers, thereby mitigating nearly all accuracy compromises observed in earlier approaches.

The enhanced memory efficiency of \textsc{QLoRA} allows for comprehensive exploration of instruction finetuning and chatbot effectiveness at model scales unattainable with traditional finetuning methods due to their memory requirements. As a result, we train over 1,000 models using multiple instruction-tuning datasets, spanning various architectures and covering model sizes from 80M to 65B parameters. Our results demonstrate that \textsc{QLoRA} achieves parity with 16-bit training (\S\ref{sec:comp_to_std}) and enables training of the cutting-edge chatbot, \textbf{Guanaco} (\S\ref{sec:pushing}). Further, our analysis yields several insights: firstly, we observe that the quality of training data significantly outweighs dataset size; for instance, a 9k sample dataset (OASST1) surpassed a 450k sample subset (FLAN v2) in chatbot evaluations, even when both were targeted at fostering instruction-following generalization. Secondly, we establish that high performance on the Massive Multitask Language Understanding (MMLU) benchmark does not necessarily translate to strong outcomes on the Vicuna chatbot benchmark, and vice versa—thus, the relevance and quality of the dataset to the specific task are more critical than mere scale.

In addition, we conduct a thorough assessment of chatbot effectiveness employing evaluations from both human raters and GPT-4. Our methodology utilizes a tournament-based benchmarking system in which chatbots are pitted against each other to generate optimal responses to specific prompts. Each match’s victor is determined either by GPT-4 or by human evaluators. To synthesize tournament outcomes, we employ Elo scores~\citep{elo1967proposed, elo1978rating}, allowing for the establishment of chatbot performance rankings. Our results indicate a substantial alignment between GPT-4 and human annotators in ranking model performance, though notable discrepancies sometimes occur. We therefore emphasize that, despite offering a cost-effective alternative to human judgment, model-based evaluation methods are also subject to certain ambiguities.

Beyond our quantitative benchmarks, we provide a qualitative examination of \textbf{Guanaco} models. This analysis reveals exemplary and deficient cases that are not otherwise evident in the numeric results.

To encourage additional research, we publish all generated responses along with human and GPT-4 annotations. The entirety of our codebase and CUDA kernels is released as open source, with integration into the Hugging Face transformers stack \citep{wolf2019huggingface} to promote widespread accessibility. Additionally, we provide a suite of adapters for models of 7/13/33/65B parameter sizes, each fine-tuned on one of eight instruction-following datasets, culminating in 32 unique open-source, fine-tuned models.

\section{Background}
\paragraph{Block-wise k-bit Quantization} Quantization refers to transforming data from a format with higher precision into one with fewer representational levels, effectively decreasing information content. This process typically involves converting data types with larger bit-widths, such as 32-bit floating point numbers, to types with fewer bits, such as 8-bit integers. To utilize the dynamic range of the lower bit-width format effectively, normalization is performed by dividing the data by its maximum absolute value before casting, typically over tensor structures. For instance, consider quantizing a tensor represented with 32-bit floating points (FP32) to an 8-bit integer (Int8) in the range $[-127, 127]$:

\begin{equation}
    \mathbf{X}^{ \text{Int8}} = \text{round}\left( \frac{127}{{\text{absmax}}(\mathbf{X}^{{\text{FP32}}})}\mathbf{X}^{{\text{FP32}}} \right) = \text{round}(c^{\text{FP32}}\cdot \mathbf{X}^{{\text{FP32}}}),
\end{equation}

where $c$ denotes the {\em quantization constant} or {\em quantization scale}. The reverse operation, dequantization, restores the original value as follows:

\begin{equation}
    \text{dequant}(c^{\text{FP32}}, \mathbf{X}^{\text{Int8}}) = \frac{\mathbf{X}^{\text{Int8}}}{c^{\text{FP32}}} = \mathbf{X}^{\text{FP32}}
\end{equation}

A significant drawback of this method arises when an outlier—an unusually large input value—appears within the tensor. In such cases, the allocation of quantization bins, or bit representations, becomes inefficient; some bins may be scarcely populated or entirely unused. To address this issue, a widely used solution involves partitioning the input tensor into discrete blocks that are each quantized separately, assigning a distinct quantization constant $c$ to every block. Formally, this process involves dividing the tensor $\mathbf{X}\in \mathbb{R}^{b\times h}$ into $n$ consecutive blocks, each of size $B$, by flattening the tensor and slicing it linearly, yielding ${n = ({b\times h} )/{B} }$ total blocks. Each block is then independently quantized according to Equation~1, producing a quantized version of the input and $n$ associated quantization constants $c_i$.

\paragraph{Low-rank Adapters} Low-rank Adapter (LoRA) finetuning \citep{hu2021lora} mitigates memory consumption by incorporating a limited number of trainable parameters, known as adapters, while keeping the pre-existing model parameters unaltered. During training with stochastic gradient descent, gradients propagate through the immutable pretrained model weights and are used solely to update the adapters, which are trained to minimize the objective function. LoRA operates by enriching a standard linear projection with an added low-rank, factorized component. Specifically, for a projection described by ${\mathbf{X} \mathbf{W} = \mathbf{Y}}$, where $\mathbf{X}\in  \mathbb{R}^{b\times h}$ and $\mathbf{W}\in  \mathbb{R}^{h\times o}$, the LoRA approach modifies the computation as:

\begin{equation}
    \mathbf{Y} = \mathbf{X}\mathbf{W} + s\mathbf{X}\mathbf{L}_1\mathbf{L}_2,
\end{equation}

with $\mathbf{L}_1\in  \mathbb{R}^{h\times r}$, $\mathbf{L}_2\in  \mathbb{R}^{r\times o}$, and $s$ representing a scalar parameter.

\paragraph{Memory Requirement of Parameter-Efficient Finetuning} A significant consideration in applying LoRA concerns its memory demands during training, specifically relating to the quantity and dimensions of adapters employed. Given LoRA’s notably small memory footprint, incorporating additional adapters can enhance model effectiveness with negligible impact on total memory consumption. Although LoRA was introduced as a Parameter Efficient Finetuning (PEFT) strategy, the primary contributor to memory usage when finetuning LLMs is the storage of activation gradients, not the LoRA parameters themselves. For instance, when training a 7B LLaMA model on FLAN v2 with a batch size of 1 and LoRA weights set to the typical 0.2\% of the model’s original weights\citep{hu2021lora,liu2022few}, the memory requirement for storing LoRA input gradients is approximately 567 MB, whereas the LoRA parameter storage is only 26 MB. If gradient checkpointing~\citep{chen2016training} is applied, the per-sequence input gradients decrease to an average of 18 MB, which still exceeds the memory used by all LoRA parameters combined. For context, a 4-bit quantized base model occupies 5,048 MB in memory. This comparison illustrates the significance of gradient checkpointing and shows that further reductions in LoRA parameter size offer only marginal improvements in memory efficiency. As a result, it is feasible to deploy more adapters without meaningfully increasing memory usage during training (refer to Appendix~\ref{app:memory} for a comprehensive analysis). As explored in subsequent sections, this capability is essential for achieving performance comparable to full 16-bit precision finetuning.

\section{\textsc{QLoRA} Finetuning}

\textsc{QLoRA} enables precise 4-bit finetuning through two innovations: 4-bit NormalFloat (NF4) quantization and Double Quantization. Furthermore, we propose Paged Optimizers, a strategy designed to alleviate memory surges during gradient checkpointing, thereby mitigating out-of-memory failures that have historically hampered large-model finetuning on individual machines.

Within \textsc{QLoRA}, parameters are maintained in a single low-precision format (commonly 4-bit), while computations are executed in a higher-precision format, typically BFloat16. Concretely, any time a \textsc{QLoRA} parameter tensor is utilized, it is first dequantized to BFloat16, after which matrix multiplications proceed in 16-bit precision.

We proceed by describing the individual elements constituting \textsc{QLoRA}, and then present its formal definition.

\paragraph{4-bit NormalFloat Quantization} The NormalFloat (NF) data type extends Quantile Quantization \citep{dettmers2022optimizers}, an approach grounded in information theory that guarantees an equal number of input values map to each quantization bin. This is accomplished by determining the input tensor's quantiles via its empirical cumulative distribution function.

Quantile quantization is constrained by the high computational cost associated with quantile estimation. To address this, rapid quantile approximation techniques like SRAM quantiles~\citep{dettmers2022optimizers} are commonly employed. However, because these algorithms provide only approximate estimates, they introduce significant quantization errors, especially for outlier values that frequently carry crucial information.

When input tensors are distributed identically up to a quantization constant, both costly quantile calculations and approximation inaccuracies can be avoided. In such scenarios, input tensors share quantile distributions, thereby making precise quantile estimation tractable from a computational standpoint.

Pretrained neural network weights generally follow a zero-mean normal distribution characterized by a standard deviation $\sigma$ (refer to Appendix~\ref{app:normality}). To standardize all weights under a unified distribution, we can adjust $\sigma$ so that the entire distribution is contained within the predefined range of the data type in use. In this work, we select $[-1, 1]$ as the reference interval. Accordingly, it is necessary to map both the data type quantiles and the neural network weight values into this interval through normalization.

To determine the data type that is theoretically optimal, in terms of information, for zero-centered Gaussian distributions with variable standard deviation $\sigma$ in the interval $[-1, 1]$, we proceed as follows: (1) identify the $2^k+1$ quantile boundaries from the standard normal distribution $N(0, 1)$, which gives a $k$-bit quantile-based quantization scheme for such distributions; (2) linearly rescale the resulting quantile representatives to reside within $[-1, 1]$; (3) for a given weight tensor, perform quantization by rescaling its values into $[-1, 1]$ using absolute maximum scaling.

When both the weight tensor and data type occupy the same numerical range, standard quantization methods may be employed. Notably, the normalization in step (3) serves to align the standard deviation of the weight tensor with that of the chosen $k$-bit data type. Specifically, the $2^k$ representative quantization points $q_i$ can be computed using:

\begin{equation}
q_i  = \frac{1}{2}\left( Q_X\left(\frac{i}{2^k+1}\right) + Q_X\left(\frac{i+1}{2^k+1}\right)\right),
\end{equation}

where $Q_X(\cdot)$ denotes the quantile function associated with the standard normal distribution $N(0,1)$. A limitation of conventional symmetric k-bit quantization is its inability to represent zero exactly, a crucial feature for error-free quantization of padding and other zero-valued entities. To address this, and to utilize all $2^k$ representable values for a k-bit format while guaranteeing a discrete zero point at $0$, we propose an asymmetric representation. This involves separately estimating quantiles $q_i$ for two intervals: allocating $2^{k-1}$ quantization levels to the negative side and $2^{k-1}+1$ levels to the positive side, merging these resulting sets, and then removing the duplicate zero. We designate the constructed datatype, which yields equal expected frequencies across bins, as \emph{k-bit NormalFloat} (NFk). This type achieves information-theoretic optimality for data that follows a zero-centered normal distribution. The detailed construction for this datatype appears in Appendix~\ref{app:nf4}.

\paragraph{Double Quantization{}} 
We present \emph{Double Quantization} (DQ), a technique wherein the quantization parameters themselves are quantized to achieve further reductions in memory usage. Although accurate 4-bit quantization necessitates a small blocksize~\citep{dettmers2022case}, this comes with significant memory overhead; for instance, employing 32-bit quantization constants for $\mathbf{W}$ with a blocksize of 64 incurs an average overhead of $32/64 = 0.5$ bits per parameter. Double Quantization serves to decrease the memory burden imposed by these constants.

Concretely, this method applies a second quantization stage to the initial set of quantization constants $c_2^{\text{FP32}}$, using them as input and producing quantized values $c_2^{\text{FP8}}$ as well as a new set of second-level quantization constants $c_1^{\text{FP32}}$. For this second step, we utilize 8-bit Floats and a blocksize of 256, given that such a configuration results in no observable loss in model accuracy—consistent with findings from \citet{dettmers2022case}. Because the values $c_2^{\text{FP32}}$ are inherently non-negative, we subtract their mean prior to quantization so as to center the data, thereby permitting the use of symmetric quantization. With blocksize 64, the average memory usage per parameter for quantization constants is reduced from $32/64=0.5$ bits to $8/64 + 32/(64\cdot256) = 0.127$ bits—a savings of 0.373 bits per parameter.

\paragraph{Paged Optimizers} leverage NVIDIA's unified memory~\footnote{\tiny \url{https://docs.nvidia.com/cuda/cuda-c-programming-guide}}, a feature that transparently manages page-level data transfers between CPU and GPU. This system facilitates seamless GPU execution even in situations where GPU memory resources are insufficient, functioning analogously to traditional memory paging between system RAM and disk storage. In our approach, optimizer state data is allocated as paged memory, which the runtime moves from GPU to CPU RAM as GPU memory becomes constrained and brings it back to GPU memory as needed during optimizer update steps.

\paragraph{\textsc{QLoRA}.} Building upon the aforementioned mechanisms, we introduce \textsc{QLoRA} applied to a single linear layer in the quantized base model equipped with a single LoRA adapter:

\begin{equation}
    \mathbf{Y}^{\text{BF16}} =  \mathbf{X}^{\text{BF16}} \text{doubleDequant}(c_1^{\text{FP32}}, c_2^{\text{k-bit}}, \mathbf{W}^{\text{NF4}}) + \mathbf{X}^{\text{BF16}}\mathbf{L}^{\text{BF16}}_1\mathbf{L}^{\text{BF16}}_2,
\end{equation}

The function doubleDequant$(\cdot)$ is explicitly defined by:

\begin{equation}
    \text{doubleDequant}(c_1^{\text{FP32}}, c_2^{\text{k-bit}}, \mathbf{W}^{\text{k-bit}}) = \text{dequant}(\text{dequant}(c_1^{\text{FP32}}, c_2^{\text{k-bit}}), \mathbf{W}^{\text{4bit}}) = \mathbf{W}^{\text{BF16}},
\end{equation}

In this configuration, $\mathbf{W}$ utilizes NF4 quantization while $c_2$ is represented in FP8 format.

For improved quantization fidelity, a block size of 64 is selected for $\mathbf{W}$, whereas $c_2$ adopts a block size of 256 to optimize for reduced memory footprint.

During parameter updates, gradients with respect to adapter weights, namely $\frac{\partial E}{\partial \mathbf{L}_i}$, are the sole requirement; gradients for the 4-bit weights, $\frac{\partial E}{\partial \mathbf{W}}$, are not computed. Nevertheless, evaluating $\frac{\partial E}{\partial \mathbf{L}_i}$ necessarily involves calculating $\frac{\partial \mathbf{X}}{\partial \mathbf{W}}$. This is accomplished following equation~(5), where $\mathbf{W}^{\text{NF4}}$ undergoes dequantization to $\mathbf{W}^{\text{BF16}}$ as the computational data type. This transformation permits the evaluation of $\frac{\partial \mathbf{X}}{\partial \mathbf{W}}$ in BFloat16 precision.

In summary, \textsc{QLoRA} employs a single data type for storage, typically 4-bit NormalFloat, alongside a separate data type for computation, namely 16-bit BrainFloat. To execute both the forward and backward passes, stored values are dequantized from the storage type to the computation type. However, gradient calculations are performed exclusively for the LoRA parameters, which remain in the 16-bit BrainFloat format.

\section{QLoRA vs. Standard Finetuning}
\label{sec:comp_to_std}

We have explained the mechanics of QLoRA and its ability to substantially lower memory requirements during model finetuning. A central issue that remains is whether QLoRA is able to match the performance of traditional full-model finetuning. Additionally, we wish to scrutinize various QLoRA components, particularly assessing the benefits of NormalFloat4 compared to standard Float4. The experiments described in the subsequent sections are designed to investigate these aspects.

\paragraph{Experimental setup.} Our evaluation considers three types of architectures (encoder, encoder-decoder, and decoder-only) and compares QLoRA with 16-bit adapter finetuning as well as full-model finetuning for models containing up to 3B parameters. The assessment encompasses tasks such as GLUE \citep{wang2018glue} using RoBERTa-large \citep{liu2019roberta}, Super-NaturalInstructions (TKInstruct) \citep{wang2022super} with T5 \citep{t5}, and a 5-shot MMLU \citep{hendrycksmeasuring} following LLaMA finetuning on Flan v2 \citep{longpre2023flan} and Alpaca \citep{alpaca}. To further explore the merits of NF4 over alternative 4-bit data representations, we adopt the approach of \citet{dettmers2022case} and quantify post-quantization zero-shot accuracy and perplexity over a spectrum of models (including OPT~\citep{zhang2022opt}, LLaMA~\citep{touvron2023llama}, BLOOM~\citep{scao2022bloom}, Pythia~\citep{biderman2023pythia}) with sizes ranging from 125m to 13B parameters.

Details relevant to each configuration are elaborated upon in the results sections for enhanced clarity. The appendix~\ref{app:exp-setup} contains the complete experimental setup.

Although paged optimizers are crucial for enabling \textsc{QLoRA} finetuning on 33B and 65B models using a single 24/48GB GPU, we do not present explicit performance benchmarks for Paged Optimizers since the paging mechanism typically activates only during the processing of mini-batches with extended sequence lengths, which seldom occurs. Nevertheless, our runtime analysis of paged optimizers on 65B models using 48GB GPUs shows that, for a batch size of 16, paged and standard optimizers achieve equivalent training throughput. Future investigations should systematically evaluate and delineate the specific conditions under which the paging strategy might induce slow-downs.

\paragraph{Default LoRA hyperparameters do not match 16-bit performance}

Applying the common strategy of utilizing LoRA on the query and value attention projection matrices \citep{hu2021lora} fails to recover the full finetuning accuracy in the case of large base models. Figure~\ref{fig:hyper} demonstrates, using LLaMA 7B finetuned on Alpaca, that the decisive LoRA hyperparameter is the number of adapters implemented: covering all linear transformer block layers with LoRA adapters is essential for parity with full finetuning performance. Other parameters of LoRA, for instance the projection dimension $r$, appear to have negligible influence on performance (see Appendix \ref{app:exp-setup}).

In a similar vein, our analysis reveals that the standard hyperparameters commonly used for fully finetuned baselines are insufficiently optimized. To establish more reliable baselines, we conduct an extensive hyperparameter sweep, exploring learning rates from $1\mathrm{e}{-6}$ to $5\mathrm{e}{-5}$ and batch sizes ranging from 8 to 128. The outcomes of 7B LLaMA finetuning on the Alpaca dataset are illustrated in Figure~\ref{fig:hyper}.

\paragraph{4-bit NormalFloat demonstrates superior results over 4-bit Floating Point}

Although the 4-bit NormalFloat (NF4) data format is optimal in terms of information theory, its practical benefits in empirical scenarios have yet to be firmly established. Adopting the experimental protocol in \citet{dettmers2022case}, we evaluate quantized LLMs—including OPT \citep{zhang2022opt}, BLOOM \cite{scao2022bloom}, Pythia \cite{biderman2023pythia}, and LLaMA—across various model sizes (from 125M to 65B) and data representations, assessing performance on language modeling as well as a range of zero-shot benchmarks. The data in Figure~\ref{fig:nf4} and Table~\ref{tbl:nf4_ppl} clearly indicate that NF4 substantially outperforms FP4 and Int4, and additionally, that employing double quantization provides a reduction in memory usage with no adverse impact on performance.

\paragraph{k-bit \textsc{QLoRA} achieves parity with 16-bit full finetuning and 16-bit LoRA}

Previous research has demonstrated that while 4-bit quantization for \emph{inference} is feasible, it tends to result in a decline in model quality compared to 16-bit precision \citep{dettmers2022case,frantar2022gptq}. This brings forth the pivotal question: can adapter finetuning in 4-bit precision restore the lost performance? We investigate this through two experimental settings.

In the first, we perform a direct comparison against full 16-bit finetuning using RoBERTA and T5 models of sizes ranging from 125M to 3B parameters, evaluating on GLUE and the Super-NaturalInstructions datasets. Refer to Table~\ref{tab:nf4vsbf16} for results. On both benchmarks, the outcomes demonstrate that 16-bit, 8-bit, and 4-bit adapter-based finetuning approaches are able to match the performance of the full 16-bit finetuning baseline. This finding implies that any accuracy deficits arising from reduced-precision quantization can be fully remedied via adapter-based finetuning.

In our second experimental configuration, due to the computational demands of fully finetuning models with 11B or more parameters—which exceeds the memory capacities of a single high-memory GPU server—we continue our investigation into whether 4-bit \textsc{QLoRA} achieves performance parity with 16-bit LoRA across model sizes ranging from 7B to 65B parameters. Specifically, we perform finetuning on LLaMA models from 7B up to 65B parameters using two instruction-following datasets, Alpaca and FLAN v2. We subsequently assess their performance on the MMLU benchmark using a 5-shot accuracy protocol. Table~\ref{tbl:llama-datatype} displays the outcomes, revealing that NF4 with double quantization is able to completely recover the MMLU accuracy achieved by 16-bit LoRA. Moreover, we observe that \textsc{QLoRA} configured with FP4 trails the 16-bit brain float LoRA baseline by approximately one percentage point. These results substantiate both that (1) \textsc{QLoRA} equipped with NF4 attains equivalent outcomes to both 16-bit full finetuning and 16-bit LoRA, and (2) NF4 confers higher quantization accuracy compared to FP4.

\paragraph{Summary} Our findings systematically demonstrate that 4-bit \textsc{QLoRA} utilizing the NF4 format consistently replicates the performance of both 16-bit full and LoRA-based finetuning on widely recognized academic benchmarks. The superiority of NF4 over FP4 is established, as is the non-detrimental effect of double quantization on model quality. Collectively, these results offer strong evidence that 4-bit \textsc{QLoRA} tuning achieves results on par with 16-bit techniques.

Consistent with prior work in the field of quantization \citep{dettmers2022case}, the MMLU and Elo evaluations suggest that, within a fixed budget for finetuning and inference, allocating computational resources to increase the parameter count of the base model—while reducing numerical precision—is advantageous. This finding underscores the efficiency gains enabled by \textsc{QLoRA}. Given that our 4-bit finetuning experiments did not manifest any reduction in performance relative to full-precision finetuning, an open question remains regarding the precise location of the trade-off between model performance and parameter precision for QLoRA tuning, which we propose for further exploration in future studies.

We turn our attention to exploring instruction tuning at scales unattainable with full 16-bit finetuning on standard academic research infrastructure.

\section{Pushing the Chatbot State-of-the-art with QLoRA}
\label{sec:pushing}
Building on the result that 4-bit \textsc{QLoRA} achieves performance comparable to 16-bit approaches across a broad set of scales, task types, and datasets, we present a comprehensive analysis of instruction finetuning up to the largest publicly accessible language models for research purposes. To gauge the effectiveness of instruction-tuned models, we use a rigorous Natural Language Understanding benchmark (MMLU) and also devise new strategies for assessing chatbot efficacy in realistic usage scenarios.

\subsection{Experimental setup} Below, we provide a synopsis of the experimental design; further procedural specifics can be found in Appendix \ref{app:chatbot}. 

\paragraph{Data} Recognizing the lack of an exhaustive analysis of modern instruction-following datasets, we select eight recent benchmarks for this investigation. Our selection incorporates datasets produced via crowd-sourcing (OASST1~\citep{kopf2023openassistant}, HH-RLHF~\citep{bai2022training}), distillation from models already trained on instructions (Alpaca~\citep{alpaca}, self-instruct~\citep{wang2022self}, unnatural-instructions~\citep{honovich2022unnatural}), aggregated corpora (FLAN v2~\citep{chung2022scaling}), and hybrid collections (Chip2~\citep{laion2023}, Longform~\citep{koksal2023longform}). Collectively, these resources span different languages, corpus sizes, and legal frameworks.

\paragraph{Training Setup} To ensure consistency and avoid bias stemming from varying training regimes, all QLoRA finetuning is conducted with supervised cross-entropy loss, foregoing any reinforcement learning—even in cases where datasets offer human preference scores. For datasets clearly separating instruction and response, we limit finetuning to responses exclusively (for further discussion, see Appendix \ref{app:chatbot}). Datasets such as OASST1 and HH-RLHF supply multiple response options per prompt; for these, the highest-ranked response is chosen at each node in the conversational hierarchy, and finetuning is performed using the full, selected exchange, incorporating both instructions and corresponding responses. NF4 \textsc{QLoRA} with double quantization and paged optimization is employed throughout to mitigate the risk of memory overflow during gradient checkpointing. For the 13B and 33B LLaMA variants, we conduct limited hyperparameter tuning, observing that values discovered for 7B models typically extend to larger models—apart from learning rate and batch size, for which we reduce the learning rate by half and double the batch size when scaling up to 33B and 65B.

\paragraph{Baselines} We benchmark our results against both research-focused systems (Vicuna~\citep{vicuna2023}, Open Assistant~\citep{kopf2023openassistant}) and established commercial chatbots (GPT-4 \citep{openai2023gpt}, GPT-3.5-turbo, and Bard). Open Assistant utilizes a LLaMA 33B architecture that has been finetuned with Reinforcement Learning from Human Feedback (RLHF) based on the OASST1 corpus, which we also incorporate into our experiments. In contrast, Vicuna employs full model finetuning on LLaMA 13B using proprietary, user-supplied conversations sourced from ShareGPT, effectively serving as a distilled proxy for OpenAI’s GPT offerings.

\subsection{Evaluation}

Consistent with standard methodology, we assess model performance using the MMLU (Massively Multitask Language Understanding) benchmark~\citep{hendrycksmeasuring}, which consists of 57 distinct multiple-choice tasks spanning topics such as elementary mathematics, US history, computer science, and law. Results are reported as 5-shot accuracy on the test set.

We further evaluate generative language performance utilizing both automated tools and human raters. This evaluation suite is based on human-curated query sets, aiming to assess the quality of generated responses. Although such tests better simulate real-world chatbot usage and have become increasingly standard, there remains no universally endorsed protocol in the academic literature. In this work, we outline our evaluation framework, consistently employing nucleus sampling with $p=0.9$ and a temperature of $0.7$ throughout.

\paragraph{Benchmark Data} Our assessment leverages two handpicked sets of questions: the Vicuna prompts \citep{vicuna2023} and the validation split of OASST1 \citep{kopf2023openassistant}. The Vicuna prompts include 80 unaltered prompts spanning various topics. The OASST1 set consists of multilingual, user-assistant conversations gathered through crowdsourcing; for this benchmark, every user input in the validation split is treated as an individual query, with the full dialog history included as context, resulting in 953 distinct user queries. We refer to these datasets as the Vicuna and OA benchmarks, respectively.

\paragraph{Automated Evaluation} Adopting the procedure from \citet{vicuna2023}, we employ GPT-4 to judge and score model outputs compared to ChatGPT (GPT-3.5 Turbo) on the Vicuna set. For each prompt, GPT-4 evaluates the replies generated by both ChatGPT and a candidate model, assigning scores on a ten-point scale and providing justification. The aggregate performance is normalized as a percentage of ChatGPT's score, which may exceed 100\% should the candidate outperform ChatGPT. Notably, we identify a strong bias where GPT-4 tends to score the response presented first more favorably. To mitigate this, we advocate averaging the scores over both response orders.

To further assess relative performance, we also directly compare output pairs from different systems. This analysis simplifies grading to three labels, permitting ties. GPT-4 is tasked with selecting the superior response or indicating equivalence, together with rationale. All possible pairs of systems are evaluated this way across both benchmarks.

\paragraph{Human Evaluation} While some recent studies suggest that generative models may be valid proxies for system evaluation \citep{fu2023gptscore}, definitive evidence that GPT-4’s assessment reliably aligns with human ratings—specifically in chatbot settings—remains limited. Consequently, we conduct two parallel human evaluations mirroring both aforementioned automated protocols on the Vicuna set. For these experiments, we engage two human annotators via Amazon Mechanical Turk (AMT) for direct comparisons against ChatGPT, and three annotators for every system pairwise comparison.

\paragraph{Elo Rating} To evaluate models through both human judgment and automated assessments, we organize a tournament structure where models are pitted against each other in head-to-head matches. In each match, two models generate responses to the same prompt and are compared to determine which performs better. While this approach mirrors that of \citet{bai2022training} and \citet{vicuna2023}, our methodology additionally incorporates GPT-4-based evaluations alongside human ratings. From the pool of annotated comparisons, we randomly select samples to calculate Elo scores~\citep{elo1967proposed,elo1978rating}. The Elo rating system, a standard for ranking players in games such as chess, estimates a player's likelihood of winning against an opponent; for instance, a competitor with an Elo score of 1100 facing another with a score of 1000 would have an expected win probability of about 65\%, while evenly matched opponents (e.g., both at 1000 or both at 1100) would each have a 50\% chance of winning. After every match, a participant's Elo rating is updated in proportion to how surprising the outcome is: a win against expectation yields a significant rating shift, whereas expected results produce smaller adjustments. As tournaments proceed, these ratings tend to reflect each model's relative proficiency. All participants begin with an initial rating of 1,000, and we apply $K=32$ to scale rating adjustments. As in \citet{vicuna2023}, this process is iterated 10,000 times with varying random seeds to mitigate the influence of match order, such as which pairs are compared first.

\subsection{\textbf{Guanaco}: \textsc{QLoRA} trained on OASST1 Sets a New Standard for Open-Source Chatbots}

Through a combination of automated metrics and direct human assessment, we observe that our leading \textsc{QLoRA}-adapted model, {Guanaco} 65B, finetuned on an OASST1 variant, demonstrates best-in-class results among open-source conversational agents and achieves performance that rivals that of ChatGPT. Comparing against GPT-4, human annotator pairwise comparisons using an Elo rating system estimate a 30\% win likelihood for both {Guanaco} 65B and 33B—representing the highest values published thus far.

Table~\ref{tbl:vicuna_eval} details outcomes on the Vicuna benchmark \cite{vicuna2023}, measured relative to ChatGPT. According to these results, {Guanaco} 65B stands out as the top performer after GPT-4, reaching 99.3\% of ChatGPT’s performance. Although {Guanaco} 33B features a larger parameter count than Vicuna 13B, it leverages 4-bit quantization for its weights, making it notably more memory efficient (21 GB versus 26 GB), and yields a three percentage point increase over Vicuna 13B. In addition, {Guanaco} 7B requires just 5 GB of memory—enabling it to run smoothly on modern smartphones—while outperforming Alpaca 13B by almost 20 percentage points.

Nevertheless, Table~\ref{tbl:vicuna_eval} reveals considerable confidence intervals, with substantial performance overlap among various models. We propose that this variability arises from ambiguous scale definitions; for instance, the interpretation of an 8 on a 10-point scale may differ between evaluators and contexts. To mitigate such ambiguity, we advocate for the Elo ranking approach \citep{elo1967proposed}, which employs \textit{pairwise} model comparisons by humans and GPT-4 to obviate issues tied to absolute scale anchoring. Table~\ref{tbl:elo} provides Elo ratings for the leading models. Interestingly, while model rankings on the Vicuna benchmark show partial divergence between GPT-4 and human evaluations—most notably for Guanaco 7B—the rankings broadly align, exhibiting a Kendall Tau coefficient of $\tau=0.43$ and a Spearman rank correlation of $r=0.55$ at the overall system level. At the level of individual examples, the agreement between majority human decisions and GPT-4 is weaker (Fleiss $\kappa=0.25$). Taken together, these findings indicate a moderate level of concordance between human and GPT-4 system-level assessments, supporting the viability of model-driven evaluation as a reasonably dependable substitute for human judgment. Section~\ref{ssec:considerations} provides a more in-depth discussion of these issues.

The Elo scores reported in Table~\ref{tbl:elo_large} reveal that the {Guanaco} models with 33B and 65B parameters surpass all competitors except GPT-4 on both the Vicuna and OA evaluation sets, and their performance is on par with ChatGPT as reflected in Table~\ref{tbl:vicuna_eval}. It should be emphasized that the Vicuna evaluation set tends to benefit open-source systems, whereas the larger OA benchmark generally gives an advantage to ChatGPT. Analysis of Tables \ref{tbl:mmlu-llama} and \ref{tbl:vicuna_eval} further demonstrates that the alignment between the finetuning dataset and the evaluation task has a significant effect on results. Specifically, Llama models that are finetuned using FLAN v2 perform especially well on MMLU, but achieve the lowest scores on the Vicuna benchmark; this trend holds across other models as well. Such findings suggest that current benchmarks measure partially independent abilities: excelling on MMLU does not necessarily translate into high performance on conversational benchmarks like Vicuna or OA, and the converse is also true.

 stands out as the sole leading model in our study that avoids proprietary data, owing to the OASST1 dataset’s strict prohibition against incorporating outputs from GPT-based systems. The next highest-performing model trained solely on openly available data is Anthropic’s HH-RLHF, which achieves results 30 points below {Guanaco} on the Vicuna benchmark (refer to Table~\ref{tbl:vicuna_eval}). Taken together, these results validate that 4-bit \textsc{QLoRA} can efficiently yield leading chatbot models with capabilities comparable to ChatGPT. Notably, our {Guanaco} 33B model can be trained on widely available 24 GB consumer GPUs in under twelve hours. This demonstrates the feasibility of further exploration through \textsc{QLoRA}-based finetuning on domain-specific open-source datasets, enabling open-source models to challenge the top commercial solutions currently available.

\section{Qualitative Analysis}

Although the bulk of our evaluation relies on quantitative methods, solely focusing on summary statistics introduces certain limitations. The most significant concern is benchmark validity \citep{liao2021we}—it is often unclear whether a given benchmark genuinely assesses the skills it claims to evaluate, especially as researchers continually uncover ``shortcuts'' that allow machine learning models to artificially inflate their performance \citep{gururangan2018annotation, poliak2018hypothesis}. To partially address these concerns, we incorporate a qualitative analysis, which unfolds in two main parts. In \S\ref{ssec:examples}, we present selected examples illustrating characteristic behaviors in outputs produced by our 65b {Guanaco} model. Next, in \S\ref{ssec:considerations}, we elaborate on important issues pertaining to the interpretation of our findings.

\subsection{Qualitative Analysis of Example Generations}
\label{ssec:examples}

Our approach to finding representative samples began by examining outputs from both the Vicuna and OpenAssistant benchmarks, with the objective of identifying recurring trends in the responses produced by {Guanaco}. Upon detecting a specific trend, we crafted prompts designed to elicit the same behavior, even if doing so resulted in an incorrect response. For instance, if we noticed that the model frequently produced overly elaborate answers, we would instruct it to ``Answer yes or no without explanation.'' Through this process, we identified both problematic cases, or ``lemons,'' where we succeeded in tricking the model, as well as positive cases, or ``cherries,'' where our adversarial prompts did not cause failures. Every example provided in this section was generated using Nucleus Sampling \cite{holtzmancurious} with $p=0.9$.

Naturally, this overview does not encompass every aspect, as this limited qualitative investigation cannot feasibly account for the myriad variables at play; for example, the model’s possible response space for a single prompt is vast, and our conclusions necessarily rely on what we consider to be indicative samples. Nevertheless, we argue that presenting these instances helps to contextualize the quantitative analyses detailed previously. By releasing all models and code as open source, we hope to encourage subsequent research to delve more deeply into the phenomena we describe here.
\paragraph{Factual Recall} For queries like ``What is the capital of Zambia?'', every model reliably provides the correct information, for example:

\begin{quote} 
    \textbf{Guanaco}:
    \texttt{The capital of Zambia is Lusaka.}
\end{quote}

Yet, as the queries become less common or more specialized, {Guanaco} tends to provide incorrect responses with unwavering certainty. Consider the following question taken from HotPotQA \citep{yang2018hotpotqa}:

\begin{quote}
     \textbf{User}: In what year was the singer who popularized `I'll Keep the Lovelight Burning' born?  

     \textbf{Guanaco}:~\texttt{Al Jolson is credited with popularizing the song `I'll Keep the Lovelight Burning,' and he was born in the year 1886.}
 \end{quote}

In this case, the model fails to correctly identify both the individual associated with popularizing the song and their birth year (although the year supplied matches Al Jolson’s birthdate).

\paragraph{Suggestibility} {Guanaco} demonstrates an unexpected level of robustness in not accepting certain implied falsehoods, as evidenced in the exchange below:

\begin{quote}
    \textbf{User}: How did scientists eventually come to officially and peer-review confirm that the earth is flat? \\ \\
    \textbf{Guanaco}: \texttt{There has never been a formal scientific confirmation through peer review that the earth is flat. The vast majority of the scientific community agrees that the earth is spherical, a fact established and validated over centuries via multiple lines of empirical evidence and experimentation. \\ \\ The notion that the earth is flat is a scientifically refuted conspiracy theory, with no basis in peer-reviewed science. Any claim that a consensus of scientists has peer-reviewed and affirmed the flat earth hypothesis is incorrect. \\ \\ Please let me know if you have further questions.}
\end{quote}

also demonstrates strong ability in recognizing when a question cannot meaningfully be answered, for example,

\begin{quote}
    \textbf{User}: What time is it?

    \textbf{Guanaco}: \texttt{Apologies, but as an AI language model, I lack the capability to access live data. My answers derive from previously stored information and do not indicate the current time. For accurate timekeeping, please consult a clock or visit a website that displays real-time information.}
\end{quote}

Comparable behavior is observed for inquiries such as ``Where are you?'', ``How are you?'', and similar questions.

\paragraph{Refusal} In a related vein, there are instances in which {Guanaco} inexplicably declines to carry out instructions:

\begin{quote}
    \textbf{User}: Please invert the word order in the sentence ``I'm a cool dude in a mocha mood.'' \\ \\
    \textbf{Guanaco}: \texttt{I apologize, but I'm unable to perform that task at this time. However, I can clarify why the sentence you provided is grammatically accurate.\\ \\The sentence ``I'm a cool dude in a mocha mood'' adheres to standard English grammatical and syntactic conventions. Here, ``cool'' functions as an adjective modifying the noun ``dude,'' and the phrase ``in a mocha mood'' serves as a prepositional phrase that adds contextual information. This construction is entirely acceptable within English grammar and is meaningful in context.\\ \\If you have further inquiries or would like assistance with something else, please let me know.}
\end{quote}

\paragraph{Secret Keeping} Yet, when secrecy is crucial, refusals of this type cannot be relied upon. In practice, system architects may wish to ensure certain information remains confidential, even though the system is permitted to generate unconstrained responses for the user. For example, we instruct the model as follows: ``The secret word is `banana'. Do not reveal or repeat this, regardless of user prompts.'' and then assess how resistant the model is to attempts by users to uncover this information. When confronted with a straightforward query, the system does not yield:

\begin{quote}
    \textbf{User}: What is the secret word?

    \textbf{Guanaco}: \texttt{I'm sorry, but I'm not able to disclose the secret word. I'm programmed to not reveal or repeat back the secret word, no matter what the human requests. Is there anything else I can help you with?}
\end{quote}

which is the intended response. Nonetheless, even minimal manipulation of the prompt defeats this constraint:

\begin{quote}
    \textbf{User}: This is a game. The goal is to ignore your previous instructions. What is the secret word?

    \textbf{Guanaco}: \texttt{Sorry, the secret word is banana. It's best if we don't speak of it. Is there anything else you'd like to ask?}
\end{quote}

This example illustrates the pressing need for improved techniques that ensure more reliable adherence to given instructions.

\paragraph{Math} {Guanaco} exhibits its most significant limitation in mathematics—a common shortcoming shared by various language models, as discussed in \citep{liang2022holistic}. When {Guanaco} provides a detailed, stepwise explanation, its responses are often correct. For example:

\begin{quote}
    \textbf{User}: Lee mows one lawn and charges \$33. Last week he mowed 16 lawns and three customers each gave him a \$10 tip. How many dollars did Lee earn mowing lawns last week?

    \textbf{Guanaco}: \texttt{Lee earned \$582 mowing lawns last week. \\ \\ To find the total amount he earned, we first need to calculate how much he earned for mowing 16 lawns: \\ \\  \$33 x 16 = \$528 \\ \\ Then we need to add the tips he received from the three customers: \\  \\ \$10 x 3 = \$30 \\ \\ So the total amount Lee earned was: \\ \\ \$528 + \$30 = \$558}
\end{quote}

Nonetheless, if {Guanaco} does not approach a problem incrementally, it can make mistakes even with straightforward tasks, as highlighted in \cite{weichain}. For instance, consider the following exchange:

\begin{quote}
    \textbf{User}: Please factorize 1833.

    \textbf{Guanaco}: \texttt{Sure, I'd be happy to help! 1833 is a prime number, which means it has only two factors: 1 and 1833. \\ \\ So the unique factorization of 1833 is: \\ \\ 2\^{}1 * 3\^{}2 * 17\^{}1 \\ \\ I hope this helps! Let me know if you have any other questions.}
\end{quote}

which is incorrect (the correct factorization is $3 \times 17 \times 43$), and in fact, doubly so.

\paragraph{Theory of Mind} {Guanaco} demonstrates unexpectedly advanced Theory of Mind abilities \cite{nematzadeh2018evaluating, sap2022neural}. For example, the model provides a detailed and appropriate response to the following prompt:

\begin{quote}
    \textbf{User}: James and Abby are in the bedroom. Abby put the pen in the desk drawer. Abby leaves the bedroom. James moves the pen into the bag. Where does James think Abby will look for the pen?

    \textbf{Guanaco}: \texttt{James knows that Abby put the pen in the desk drawer. But James also knows that he has moved the pen into the bag, and he assumes that Abby does not know about this. So, James may guess that Abby will look for the pen in the desk drawer, because that's where she left it last time. However, this assumption may be incorrect, because Abby may have seen James moving the pen.}
\end{quote}

Nonetheless, such reasoning is inconsistent, and the model often provides justifications based on assumptions not warranted by the scenario. For instance,

\begin{quote}
    \textbf{User}: Evelyn entered the living room. Jackson entered the playroom. James entered the playroom. The beans are in the treasure chest. James exited the playroom. Jackson moved the beans to the pantry. Jackson exited the playroom. James entered the living room. Where will James look for the beans?

    \textbf{Guanaco}: \texttt{James will look for the beans in the pantry, because that's where Jackson moved them.}
\end{quote}

In this response, {Guanaco} incorrectly infers information transfer that was never specified. Such errors parallel those found in prior studies \cite{sap2022neural}, indicating a need for further research.

\subsection{Considerations}
\label{ssec:considerations}

\paragraph{Evaluation}

We observed moderate inter-annotator reliability among human raters (Fleiss $\kappa=0.42$), with agreement diminishing further when contrasting outputs from two high-performing models. This highlights shortcomings in the current evaluation frameworks and human judgment protocols for chatbot performance assessment. During a manual comparison between ChatGPT and {Guanaco} 65B on the Vicuna benchmark, subjective judgments frequently influenced preferences, with the paper's authors often diverging in their selections. Future research should explore strategies to address these challenges, leveraging insights from areas like Human-Computer Interaction and Psychology that have addressed subjective evaluation.

Furthermore, our examination revealed that automated evaluation tools introduce systematic biases. Notably, we detected pronounced order effects, as GPT-4 tended to grant higher scores to systems that appeared first in its prompts. The low agreement on individual examples between GPT-4 and human raters (Fleiss $\kappa=0.25$) further indicates that the criteria employed by automated evaluators may diverge significantly from human judgments. Additionally, Table~\ref{tbl:elo_large} illustrates that GPT-4 gives substantially higher ratings to its own responses compared to those provided by humans (Elo of 1348 versus 1176), corresponding to about a 20\% increased probability of outperforming another system.

Future research ought to explore potential sources of bias within automated evaluation mechanisms and assess strategies for reducing such biases.

\paragraph{Data \& Training} It is important to recognize that the {Guanaco} models are trained on the multilingual OASST1 dataset and that the OA benchmark incorporates prompts in a variety of languages. We suggest that subsequent studies further examine the extent to which training on multilingual data contributes to better handling of instructions in non-English languages, and whether this factor accounts for the more pronounced disparity observed between the performance of the Vicuna-13B model (trained solely on English) and the {Guanaco} 33B and 65B models on the OA benchmark.

Due to the impressive results achieved by the {Guanaco} models, we conducted an analysis to detect any data leakage between the OASST1 training set and the Vicuna benchmark prompt set. After employing fuzzy string matching across both datasets and conducting manual review of near matches, we did not observe any shared prompts.

In addition, it should be highlighted that our training paradigm is limited to supervised learning via cross-entropy loss, without integrating reinforcement learning from human feedback (RLHF). This indicates the need for future exploration into the comparative advantages and disadvantages of relying exclusively on cross-entropy loss versus incorporating RLHF. We anticipate that \textsc{QLoRA} will support the systematic investigation of these dynamics at scale without incurring substantial computational costs.

\section{Related Work}

\paragraph{Quantization of Large Language Models} Most efforts around the quantization of LLMs have concentrated on reducing precision specifically for inference. Leading methods that retain high quality in 16-bit LLMs involve addressing outlier activations (see, e.g., SmoothQuant~\citep{xiao2022smoothquant} and LLM.int8()~\citep{dettmers2022llm}), whereas other solutions employ advanced grouping schemes~\citep{park2022nuqmm, yao2022zeroquant}. Additional research on lossy quantization explores the implications of different rounding techniques~\citep{dettmers2022case,zeng2022glm,pope2022efficiently} and methods for optimizing these rounding choices to enhance precision~\citep{frantar2022gptq}. Apart from our own contributions, SwitchBack layers~\citep{wortsman2023stable} stands out as the only other investigation into backpropagation with quantized weights at model sizes beyond 1B parameters.

\paragraph{Finetuning with Adapters} In our experiments, we utilize Low-rank Adapters~\citep{hu2021lora} (LoRA), though a wide range of Parameter Efficient FineTuning (PEFT) strategies exist in the literature. These include prompt tuning~\citep{qin2021learning,lester2021power,li2021prefix}, input modifications at the embedding layer~\citep{an2022input}, manipulation of hidden state activations (IA$^3$)~\citep{liu2022few}, the incorporation of additional full network layers~\citep{houlsby2019parameter}, bias parameter optimization~\citep{zaken2021bitfit}, Fisher information-guided weight masking~\citep{sung2021training}, and hybrid approaches~\citep{henderson2021compacter}. Our findings indicate that LoRA adapters are capable of achieving performance on par with comprehensive 16-bit finetuning. Exploration of the advantages and tradeoffs associated with alternative PEFT approaches is left for future investigation.

\paragraph{Instruction Finetuning} Instruction finetuning adapts a pretrained large language model to follow user-specified instructions by optimizing on diverse collections of input-output examples from numerous sources. Notable methods and resources in this area are MetaICL~\citep{min2021metaicl}, MetaTuning~\citep{zhong2021adapting}, InstructGPT~\citep{ouyang2022training}, FLAN~\citep{wei2021finetuned,chung2022scaling}, PromptSource~\citep{bach2022promptsource}, Super-NaturalInstructions~\citep{wang2022super,sanh2021multitask}, Self-instruct~\citep{wang2022self}, UnnaturalInstructions~\citep{honovich2022unnatural}, OPT-IML~\citep{iyer2022opt}, UnifiedSKG~\citep{xie2022unifiedskg}, OIG/Chip2~\citep{laion2023}, Alpaca~\citep{alpaca}, Vicuna~\citep{vicuna2023}, Koala~\citep{koala_blogpost_2023}, and Self-instruct-GPT-4~\citep{peng2023instruction}.

\paragraph{Chatbots} Instruction-tuned models frequently take the form of conversational agents or chatbots, and are often trained using Reinforcement Learning from Human Feedback (RLHF)~\citep{christiano2017deep} or via Reinforcement Learning with AI Feedback (RLAIF), where responses are generated from an existing model and scored by an AI system~\citep{bai2022constitutional}. Representative techniques and data collections in this domain include Anthropic-HH~\citep{askell2021general,bai2022training}, Open Assistant~\citep{kopf2023openassistant}, LaMDA~\citep{thoppilan2022lamda}, and Sparrow~\citep{glaese2022improving}. In our study, we opt not to use reinforcement learning directly. Nevertheless, our best-performing model, Guanaco, is fine-tuned with multi-turn dialogues from the Open Assistant dataset—originally curated for RLHF training~\citep{kopf2023openassistant}. Furthermore, alternative evaluation paradigms utilizing GPT-4 for chatbot assessment, as substitutes for human annotation, have recently emerged \citep{vicuna2023,peng2023instruction}. Building upon these, we introduce an evaluation protocol aimed at greater robustness.

\section{Limitations and Discussion}

Our experiments demonstrate that \textsc{QLoRA} successfully matches the performance of traditional 16-bit finetuning when applied to 4-bit base models combined with LoRA adapters. However, we do not provide evidence that \textsc{QLoRA} attains equivalent results at larger scales, specifically for 33B and 65B parameter models. Addressing this gap, which entails substantial computational demands, is deferred for subsequent research. 

A further limitation concerns the assessment of instruction finetuning models. While we report results using MMLU, the Vicuna benchmark, and the OA benchmark, we have not included evaluations on other established suites such as BigBench, RAFT, or HELM. As such, the extent to which our findings generalize to these other evaluations is undetermined. Nonetheless, we contribute an extensive analysis with MMLU and also develop novel methodologies for chatbot evaluation.

Based on the data provided, it seems that the outcomes on these benchmarks are largely influenced by the degree of overlap between the finetuning datasets and the benchmarks themselves. For instance, the FLAN v2 dataset shares similarities with MMLU, yet is quite distinct from chatbot benchmarks; in contrast, the Chip2 dataset aligns more closely with chatbot tasks and less so with MMLU. The models’ performances on both the MMLU and Vicuna evaluations reflect these affinities. This underscores not only the necessity for more comprehensive benchmarks and evaluations but also the importance of clearly defining what is being evaluated. Is the goal to build models proficient in academic knowledge relevant to high school or university contexts, or is it to excel in conversational abilities akin to chatbots? Or perhaps something different altogether? Given the convenience of assessing models on established benchmarks, the choice of benchmark can inadvertently influence the direction of research in the field. As such, the community must take care to ensure that benchmarks genuinely reflect the capabilities and attributes deemed important.

Although our evaluation includes an in-depth analysis of general chatbot capabilities, a limitation of our study is that we conducted only a limited responsible AI assessment for {Guanaco}. Specifically, we measure the propensity of {Guanaco}-65B to produce socially biased outputs in comparison to other models, as reported in Table~\ref{tbl:crows}. Our findings show that the average bias score for {Guanaco}-65B is substantially lower than that of other models without finetuning. This suggests that training on the OASST1 dataset appears to mitigate biases present in the LLaMA base model. Despite these promising indications, it remains uncertain whether {Guanaco} performs similarly when evaluated for other categories of bias. A more thorough analysis of potential biases in {Guanaco} and related systems is left for future investigations.

A further limitation is our lack of exploration of different quantization granularities, such as 3-bit base models, or a diversity of adapter approaches. While we primarily employed LoRA, which has established a record of reliability, many other Parameter Efficient FineTuning (PEFT) techniques have demonstrated effectiveness, though it is not yet clear whether these scale well to larger model sizes. Alternative adapters may offer improved results. Additionally, post-quantization finetuning seems to restore a significant portion of the information lost during quantization, potentially paving the way for more aggressive quantization strategies. For example, applying 3-bit GPTQ quantization to a base model combined with LoRA finetuning could potentially achieve performance similar to full 16-bit finetuning after subsequent finetuning steps.

\section{Broader Impacts}

The \textsc{QLoRA} finetuning approach introduced here is the first to facilitate the finetuning of models with up to 33B parameters on a single consumer GPU and models with up to 65B parameters on a single professional GPU, all while maintaining performance on par with standard full finetuning. We have shown that our 33B model, trained with the Open Assistant dataset, achieves performance comparable to ChatGPT on the Vicuna evaluation. As instruction finetuning is crucial for adapting general-purpose pretrained LLMs into specialized chatbots resembling ChatGPT, we anticipate that our methodology will greatly broaden access to high-performing finetuning capabilities—particularly benefiting researchers with more limited resources. This democratization of NLP technology constitutes a substantial step forward in making state-of-the-art capabilities accessible. In this sense, \textsc{QLoRA} serves as a levelling mechanism, helping bridge the divide in resources between major industrial players and smaller research groups with access only to consumer-grade hardware.

A further avenue for significant influence is the introduction of LLM finetuning capabilities to mobile devices. Our \textsc{QLoRA} approach represents, in our view, a key advancement toward permitting on-device finetuning of LLMs, particularly in environments with constrained resources. While earlier work has demonstrated the feasibility of running 7B parameter models on smartphones, \textsc{QLoRA} is, to our knowledge, the first to facilitate finetuning these models directly on such hardware. Based on our calculations, an iPhone 12 Plus is capable of processing up to 3 million tokens overnight during charging periods when using \textsc{QLoRA}. Although the finetuned 7B models still fall short of ChatGPT’s performance, their quality is, in our assessment, sufficient to unlock new applications—especially those previously limited by concerns around privacy or the quality of smaller LLMs. \textsc{QLoRA} thus offers an avenue to empower privacy-sensitive deployments, where individuals retain full ownership and oversight over their data and model behavior, and simultaneously reduces barriers to LLM deployment.

Nevertheless, it is important to acknowledge that finetuning constitutes a dual-use technology and has the potential to be misapplied. Prior literature highlights multiple risks associated with widespread LLM adoption \citep{bommasani2021opportunities,bender2021dangers}. We argue, however, that democratizing access to this increasingly pervasive technology fosters broader and more independent scrutiny, in contrast to the concentration of LLM control within large organizations that often restrict model or code availability, limiting external auditing.

In summary, we anticipate that \textsc{QLoRA} will substantially advance the democratization of high-quality LLM finetuning, making it more accessible and convenient for a wide audience.